% SPDX-License-Identifier: CC-BY-SA-4.0
% Copyright (c) 2026 Rafael Beraldo
% Parte 6: como se tornar um radioamador
% Incluído por presentation.tex; compile a partir dele.

% =============================================================================
\section{Como se tornar um radioamador}
% =============================================================================

% Fluxograma: figuras/licenca.tex (3 passos, um por clique). Dados: gov.br/anatel, Res. 777/2025
\begin{frame}{O caminho das pedras: três passos na Anatel}
  \centering
  \figurapassos{licenca}{3}

  \smallskip
  \raggedright\small
  Só o passo~1 é obrigatório: o \alert{COER} é a sua ``carteira de motorista''; o indicativo é a
  ``placa'' da estação e vem no passo~3 \citep{anatel2026}.
\end{frame}

\begin{frame}{Prova, classes e custos}
  \small
  \begin{columns}[T,onlytextwidth]
    \column{0.5\textwidth}
      \textbf{A prova (habilitação)}
      \begin{itemize}
      \item Marcada pelo sistema SEC da Anatel
      \item Comece pela \alert{classe C}: técnica e ética, legislação e noções de eletrônica;
        51\,\% em cada matéria
      \item Apostilas e simulados no site da LABRE
      \end{itemize}
    \column{0.46\textwidth}
      \textbf{Quanto custa (estação móvel)}
      \begin{itemize}
      \item Prova e COER: \alert{grátis}
      \item Outorga: R\$~20, uma vez
      \item Licença: R\$~37 por 10 anos
      \item Por ano: R\$~10, até 31 de março
      \end{itemize}
  \end{columns}

  \smallskip
  \begin{atencao}
    Muitos guias na internet ainda descrevem as regras antigas (prova presencial, taxa da LABRE).
    Desde a Resolução 777/2025 a prova é online e gratuita: confira em
    \linkexterno{\linkAnatelRadioamador}{gov.br/anatel}.
  \end{atencao}
\end{frame}

%%% Local Variables:
%%% TeX-master: "../presentation"
%%% fill-column: 100
%%% End:
